\documentclass[11pt]{article}
\usepackage[margin=1in]{geometry}
\usepackage{amsmath,amssymb}
\usepackage{booktabs}
\usepackage{tabularx}
\usepackage{enumitem}
\usepackage{xcolor}
\usepackage[hidelinks]{hyperref}
\setlist{itemsep=2pt,topsep=4pt}

\definecolor{accent}{HTML}{B4540A}
\newcommand{\say}[1]{\begin{quote}\itshape #1\end{quote}}
\newcommand{\q}[1]{\par\medskip\noindent\textbf{\color{accent}Q: #1}\par\noindent}
\newcommand{\key}[1]{\par\smallskip\noindent\fbox{\parbox{\dimexpr\linewidth-2\fboxsep-2\fboxrule}{\textbf{Key point.} #1}}\par\smallskip}

\title{Project 0 Pitch: Presentation Prep\\[4pt]
\large Which edges can a GNN trust? A Bayesian view of edge reliability}
\author{Luis Mendez \quad CSE5825 Bayesian Machine Learning}
\date{}

\begin{document}
\maketitle

\noindent\textbf{How to use this sheet.} Sections~1--2 are what you \emph{say}: the
3-minute plan and a slide-by-slide script (say it in your own words, don't memorize).
Sections~3--9 are what you need to \emph{know} so you can answer anything: background, the
full model, a worked example, metrics, the experiment plan, related work and links to the
course. Sections~10--12 prepare you for questions. This is rehearsal material, not the
written Project~0 submission, which you must write yourself.

\setcounter{tocdepth}{1}
\tableofcontents
\newpage

%================================================================
\section{The pitch in one breath}

\say{GNNs trust every edge, but edges can be noisy or attacked. I treat whether each edge is
real as a hidden variable and infer it with Bayes' rule: an LLM reads the two nodes' texts to
give a prior, the graph structure gives a likelihood, and the GNN weights each message by the
posterior. Because I inject the fake edges myself, I can check every posterior against
ground truth.}

\subsection*{Three things the audience must remember}
\begin{enumerate}
  \item \textbf{The problem:} GNNs blindly trust edges, and fake edges can be made to look
        plausible.
  \item \textbf{The idea:} posterior $\propto$ LLM text prior $\times$ structural likelihood,
        for every edge.
  \item \textbf{Why it is credible:} injected edges give ground truth, so every claim is
        checkable.
\end{enumerate}

\subsection*{Time budget (about 3 minutes)}
\begin{center}
\begin{tabular}{clr}
\toprule
Slide & Content & Time \\
\midrule
1 & Title and question & 15\,s \\
2 & Motivation & 30\,s \\
3 & Datasets & 30\,s \\
4 & Approach (Bayes' rule for every edge) & 40\,s \\
5 & Three hypotheses & 35\,s \\
6 & Why this is testable & 15\,s \\
7 & How I'll check it, and the ask & 15\,s \\
\midrule
 & \textbf{Total} & \textbf{3:00} \\
\bottomrule
\end{tabular}
\end{center}
If you run long, cut slide~6 to one sentence; slide~5 already implies it. If you run short,
spend the extra time on slide~4.

%================================================================
\section{Slide-by-slide script}

\subsection*{Slide 1: Title (15\,s)}
\say{Hi, I'm Luis. My project asks a simple question about graph neural networks: which edges
in a graph should the model actually trust? I want to answer it with a Bayesian model that
uses an LLM to read the node texts.}

\subsection*{Slide 2: Motivation (30\,s)}
\begin{itemize}
  \item GNNs pass messages along edges and assume every neighbor is reliable evidence.
  \item Edges can be noisy or deliberately injected. An attacker can connect nodes whose texts
        \emph{look} related, so the fake edges are hard to spot.
  \item Two kinds of evidence exist (structure and text), but current defenses combine them
        with hand-made rules and give no uncertainty.
  \item \textbf{Goal:} treat each edge's reliability as a hidden variable and infer it.
\end{itemize}
\emph{Transition:} ``To study this I need graphs with text, and a way to know which edges are
fake.''

\subsection*{Slide 3: Datasets (30\,s)}
\begin{itemize}
  \item Four standard text-attributed graphs: Cora, CiteSeer, PubMed (paper citations) and
        WikiCS (Wikipedia CS articles).
  \item Each node has raw text, a topic label, and citation or hyperlink edges.
  \item I add fake edges between textually similar non-neighbors (5--40\% of the edge count)
        and an LLM score for every edge, with SBERT similarity as a cheaper baseline.
  \item Punchline: I know exactly which edges are fake.
\end{itemize}
\emph{Transition:} ``Here is how the model decides which edges to trust.''

\subsection*{Slide 4: Approach (40\,s), the most important slide}
\[
\underbrace{p(z_{ij}=1 \mid s_{ij}, d_{ij})}_{\text{posterior}}
\;\propto\;
\underbrace{p(z_{ij}=1 \mid s_{ij})}_{\text{prior from LLM text score}}
\;\times\;
\underbrace{p(d_{ij} \mid z_{ij}=1)}_{\text{likelihood from structure}}
\]
Walk left to right: \textbf{prior card} (the LLM reads both texts), \textbf{likelihood card}
(do the endpoints' representations agree?), \textbf{posterior card} (the probability the edge
is real). Then point at the bottom strip: messages weighted by the posterior, a posterior on
the GNN weights too, predictions with uncertainty.
\say{In one sentence: posterior equals text prior times structural likelihood, and the GNN
uses the posterior as a soft trust weight.}

\subsection*{Slide 5: Hypotheses (35\,s)}
\begin{enumerate}[label=\textbf{H\arabic*}]
  \item \textbf{Text and structure together find fake edges.} \emph{Test:} AUROC for
        fake-edge detection.
  \item \textbf{Uncertainty flags attacked nodes.} \emph{Test:} ECE and predictive
        uncertainty vs.\ attack size.
  \item \textbf{LLMs make a better prior than embeddings.} \emph{Test:} AUROC with LLM vs.\
        SBERT prior, per dataset.
\end{enumerate}

\subsection*{Slide 6: Why this is testable (15\,s)}
\say{I inject the fake edges myself, so I know the true answer for every edge and can check
the posteriors directly, at every attack size, not just final accuracy.}

\subsection*{Slide 7: Checks and ask (15\,s)}
\say{To trust the results I'll first test on simulated graphs, then compare evidence sources,
check calibration, and compare three posterior approximations. I'm looking for one or two
teammates interested in graph learning, LLMs or Bayesian modeling. Thanks!}

%================================================================
\section{Background you should know}

\subsection{Graph neural networks}
A graph $G=(V,E)$ has node features $X$ (here, encoded node text) and adjacency matrix $A$.
A graph convolutional network (GCN) layer computes
\[
H^{(\ell+1)} = \sigma\!\left(\tilde D^{-1/2}\tilde A\,\tilde D^{-1/2} H^{(\ell)} W^{(\ell)}\right),
\qquad \tilde A = A + I,
\]
so each node averages its neighbors' (degree-normalized) representations and applies a
learned linear map $W^{(\ell)}$. Two layers mean information travels two hops.
\textbf{Task:} semi-supervised node classification. A few nodes are labeled, and the model
predicts the rest (transductive setting).

\subsection{Why fake edges hurt}
GNNs work because of \textbf{homophily}: linked nodes usually share a label. A fake edge
between nodes of different classes pulls a node's representation toward the wrong class.
Because messages travel several hops, a few edges can flip many predictions.

\subsection{The attack setting}
\begin{itemize}
  \item \textbf{Candidate pool:} for each node, rank other nodes' texts with BM25 (a classic
        keyword-matching relevance score from information retrieval) and keep the top
        non-neighbors. These pairs look textually plausible.
  \item \textbf{Selection:} a white-box projected gradient descent (PGD) attack relaxes ``add
        this edge?'' to a continuous variable, follows the gradient of the classification
        loss, then keeps the $B$ most damaging edges.
  \item \textbf{Budget:} $B$ = 5, 10, 20, 30 or 40\% of the original edge count; edges are
        added at evaluation time only. The attacker cannot change node texts.
\end{itemize}
\key{This is the hard case for text: fake edges are chosen to look textually reasonable, so
the structural signal has to carry part of the load. That is why combining both matters.}

\subsection{Text encoders and LLM scoring}
\textbf{SBERT} (Sentence-BERT) maps each text to a vector; similarity is the cosine between
two vectors, computed independently for each text. An \textbf{LLM judge} reads both texts
together in one prompt and outputs a rating of whether the link makes sense, so it can use
the relation between the texts, not just topic overlap. It is slower and its scores are not
automatically calibrated.

%================================================================
\section{The full model}

\subsection{Variables}
\begin{center}
\begin{tabular}{lll}
\toprule
Symbol & Meaning & Hidden or observed \\
\midrule
$z_{ij}\in\{0,1\}$ & edge $(i,j)$ is real (1) or injected (0) & hidden \\
$W$ & GNN weights & hidden \\
$s_{ij}$ & LLM (or SBERT) score for edge $(i,j)$ & observed \\
$d_{ij}$ & distance between endpoint representations & computed \\
$X$, $E$ & node texts and observed edge list (incl.\ injected) & observed \\
$y_{\text{train}}$ & labels of training nodes & observed \\
\bottomrule
\end{tabular}
\end{center}

\subsection{Prior: calibrated LLM score}
\[
p(z_{ij}=1 \mid s_{ij}) = \sigma(\alpha + \beta\, s_{ij}),
\]
with $\alpha,\beta$ learned (or given a prior). This turns a raw LLM rating into a calibrated
probability: if the LLM is overconfident, $\beta$ shrinks.
\emph{Extension:} share $\alpha,\beta$ across datasets through a hierarchical prior.

\subsection{Likelihood: structural consistency}
After propagation, compare endpoint representations, e.g.
$d_{ij} = \bigl\lVert h_i/\sqrt{\deg_i} - h_j/\sqrt{\deg_j}\bigr\rVert$. Model it as a
two-component mixture:
\[
d_{ij} \mid z_{ij}=1 \sim \text{Gamma}(a_1,b_1)\ \text{(small, tight)},\qquad
d_{ij} \mid z_{ij}=0 \sim \text{Gamma}(a_0,b_0)\ \text{(larger, spread out)}.
\]

\subsection{Posterior for one edge}
\[
q_{ij} \equiv p(z_{ij}=1 \mid s_{ij}, d_{ij}) =
\frac{\pi_{ij}\, p(d_{ij}\mid z=1)}{\pi_{ij}\, p(d_{ij}\mid z=1) + (1-\pi_{ij})\, p(d_{ij}\mid z=0)},
\qquad \pi_{ij} = \sigma(\alpha+\beta s_{ij}).
\]

\subsection{Inference loop (EM-style)}
\begin{enumerate}
  \item Initialize $q_{ij} = \pi_{ij}$ (trust the prior).
  \item Propagate with messages weighted by $q_{ij}$ to get representations $h$.
  \item Compute $d_{ij}$ for every edge.
  \item \textbf{E-step:} update $q_{ij}$ with the formula above.
        \textbf{M-step:} refit $a_0,b_0,a_1,b_1$ (and $\alpha,\beta$).
  \item Repeat until the $q_{ij}$ stop changing; train $W$ on the training labels.
\end{enumerate}
Alternatives worth naming: mean-field variational inference over all $z_{ij}$ jointly, or
Gibbs sampling of $z$ (MCMC) for small graphs like Cora.

\subsection{Uncertainty over the GNN weights}
Approximate $p(W\mid\mathcal D)$ with MC dropout, deep ensembles or mean-field VI, then
average predictions:
$p(y_v \mid \mathcal D) \approx \frac1T\sum_{t=1}^T p(y_v \mid X, E, q, W_t)$.
Predictive uncertainty splits into two parts:
\[
\underbrace{\mathbb H\bigl[\bar p(y_v)\bigr]}_{\text{total}}
= \underbrace{\mathbb E_{W}\,\mathbb H\bigl[p(y_v\mid W)\bigr]}_{\text{aleatoric (data noise)}}
+ \underbrace{\mathbb I(y_v; W)}_{\text{epistemic (model uncertainty)}} .
\]
H2 predicts the \emph{epistemic} part rises on attacked nodes.

%================================================================
\section{Worked example: why combining evidence helps}
A fake edge was chosen to look textually plausible, so the LLM is fooled:
$\pi_{ij} = p(z=1\mid s) = 0.8$. But the endpoints' representations disagree, giving a large
$d_{ij}$ with $p(d\mid z=1) = 0.1$ and $p(d\mid z=0) = 0.6$:
\[
q_{ij} = \frac{0.8\times0.1}{0.8\times0.1 + 0.2\times0.6} = \frac{0.08}{0.20} = 0.40 .
\]
Text alone said 80\% real; with structure the posterior drops to 40\%, so the message is
down-weighted by more than half. Conversely, a real edge between two nodes whose
representations happen to differ is protected by a strong text prior.
\key{Each source covers the other's blind spot: text catches structurally plausible fakes,
structure catches textually plausible fakes. This example is your best answer to ``why
both?''}

%================================================================
\section{Metrics in detail}
\begin{description}[style=nextline]
  \item[AUROC (H1, H3)] Rank all edges by $1-q_{ij}$. AUROC is the probability that a random
    fake edge ranks above a random real edge. 0.5 is chance, 1.0 is perfect. It does not
    depend on a threshold or on the fake/real ratio.
  \item[AUPRC] Area under precision--recall; more informative when fake edges are rare (at a
    5\% budget only about 1 in 20 edges is fake).
  \item[Expected calibration error (H2)] Bin predictions by confidence into $B$ bins:
    $\text{ECE} = \sum_{b=1}^{B} \frac{|B_b|}{n}\,\bigl|\text{acc}(B_b) - \text{conf}(B_b)\bigr|$.
    Show it with a reliability diagram (confidence vs.\ accuracy; perfect is the diagonal).
  \item[Negative log-likelihood and Brier score] Proper scoring rules that reward both
    accuracy and honest confidence.
  \item[Uncertainty AUROC (H2)] How well predictive uncertainty separates attacked nodes
    (those touching a fake edge) from clean ones.
  \item[Accuracy under attack] Node-classification accuracy at each budget, to show the
    method still classifies well.
\end{description}

%================================================================
\section{Experiment plan}
\begin{center}
\small
\begin{tabularx}{\linewidth}{>{\raggedright\arraybackslash}p{1.1cm}>{\raggedright\arraybackslash}X>{\raggedright\arraybackslash}X>{\raggedright\arraybackslash}X}
\toprule
& Compare & Metric & Supported if \\
\midrule
H1 & LLM prior only; structure only; combined posterior & AUROC/AUPRC for fake edges, per budget & combined beats both on most datasets and budgets \\
H2 & standard GNN vs.\ Bayesian GNN (3 approximations) & ECE, NLL, uncertainty AUROC vs.\ budget & Bayesian ECE grows more slowly; uncertainty flags attacked nodes \\
H3 & LLM prior vs.\ SBERT prior (and BM25) & fake-edge AUROC, accuracy & LLM wins, and the gain tracks how well text predicts labels \\
\bottomrule
\end{tabularx}
\end{center}
\textbf{Baselines to mention:} GCN, GAT, a structure-only robust GNN (e.g.\ GNNGuard or
RUNG-style robust aggregation), and the same model with a uniform prior.
\textbf{Repetitions:} several random seeds and splits; report mean $\pm$ standard deviation.
\textbf{Sanity check first:} contextual SBM graphs with planted fake edges.

%================================================================
\section{How it connects to the course}
\begin{itemize}
  \item \textbf{Latent variable model:} $z_{ij}$ is a discrete latent variable per edge;
        the likelihood is a mixture model.
  \item \textbf{EM:} the inference loop is EM (E-step on $z$, M-step on mixture and prior
        parameters).
  \item \textbf{Variational inference:} mean-field $q(z)$ and Bayes-by-backprop for $W$; the
        ELBO is the objective.
  \item \textbf{MCMC:} Gibbs sampling of $z$ on small graphs as a gold standard to compare
        VI against.
  \item \textbf{Hierarchical modeling:} sharing prior parameters across the four datasets.
  \item \textbf{Model checking:} posterior predictive checks and calibration.
  \item \textbf{Graphical models:} the graph itself is uncertain, which is a model over
        structure, not just over node labels.
\end{itemize}

%================================================================
\section{Related work (verify each before citing)}
\begin{itemize}
  \item \textbf{GNN basics:} GCN (Kipf \& Welling, ICLR 2017); GAT (Veli\v{c}kovi\'c et al.,
        ICLR 2018); GraphSAGE (Hamilton et al., NeurIPS 2017).
  \item \textbf{Attacks:} Nettack (Z\"ugner et al., KDD 2018); PGD topology attack (Xu et
        al., IJCAI 2019).
  \item \textbf{Defenses:} GNNGuard (Zhang \& Zitnik, NeurIPS 2020); RGCN (Zhu et al., KDD
        2019); Pro-GNN (Jin et al., KDD 2020); GCN-Jaccard (Wu et al., IJCAI 2019).
  \item \textbf{Bayesian/uncertain graphs:} Bayesian GCN with a random graph (Zhang et al.,
        AAAI 2019); Graph Posterior Network (Stadler et al., NeurIPS 2021).
  \item \textbf{Calibration of GNNs:} Wang et al., ``Be Confident!'' (NeurIPS 2021).
  \item \textbf{Uncertainty in deep nets:} MC dropout (Gal \& Ghahramani, ICML 2016); deep
        ensembles (Lakshminarayanan et al., NeurIPS 2017).
  \item \textbf{Text and LLMs on graphs:} SBERT (Reimers \& Gurevych, EMNLP 2019); LLMs for
        text-attributed graphs, e.g.\ TAPE (He et al., ICLR 2024).
  \item \textbf{Datasets:} Sen et al.\ (2008) for Cora/CiteSeer/PubMed; Mernyei \& Cangea
        (2020) for WikiCS.
\end{itemize}

%================================================================
\section{Likely questions and answers}

\q{Why make it Bayesian instead of using an existing robust GNN?}
Existing defenses output a weight per edge with no uncertainty and combine text and structure
by hand. The Bayesian version gives a probability per edge that I can check against ground
truth, combines both sources in a principled way, and yields calibrated predictions.

\q{What is new here?}
Treating edge trust as a posterior with an \emph{LLM-based prior} and a \emph{structural
likelihood}, and evaluating the posterior itself (detection and calibration), not just
accuracy.

\q{Isn't $d_{ij}$ circular? Representations depend on the edges you are judging.}
Yes, so inference alternates, exactly like EM: representations, then posteriors, then
representations again. A simpler fallback computes $d_{ij}$ from a text-only encoder, which
breaks the loop.

\q{Does the EM loop converge?}
EM never decreases the likelihood when the representation step is held fixed; with the
alternating GNN step it is a heuristic, so I would monitor $q$ for stability and report it.

\q{You assume edges are independent given the data. Is that reasonable?}
It is a simplification (mean-field). Attackers often target one node with several edges, so
edges are correlated. A natural extension is a node-level prior on how many edges are fake.

\q{Why would an LLM beat SBERT?}
SBERT compares two separately computed embeddings; an LLM reads both texts together. It may
not win. That is what H3 tests, and a negative result is still useful.

\q{Isn't scoring every edge with an LLM expensive?}
Cora and CiteSeer have only thousands of edges. PubMed and WikiCS are larger, so I would use
a small open model with cached scores, score a subset, or fall back to SBERT there.

\q{The attacker picks textually similar pairs. Doesn't that defeat the text prior?}
That is deliberate: it is the hard case for text, so the structural likelihood must do real
work. An attacker who adapts to the LLM itself is future work.

\q{How will you know the LLM's scores are trustworthy?}
I do not assume they are. The logistic prior learns how much to trust them, and I check
calibration directly: among edges the LLM calls plausible, how many are real?

\q{Why not just threshold the posterior and delete edges?}
A soft weight keeps information when the model is unsure and avoids deleting real edges; a
threshold also has to be tuned. I can compare both.

\q{Are the ``real'' edges really all real?}
No. The original graphs have their own noise. Ground truth covers injected edges only, which
I state as a limitation.

\q{Why MC dropout, ensembles \emph{and} VI?}
They are the standard cheap approximations and are known to differ in calibration; comparing
them is part of the Bayesian question.

\q{Why not full MCMC?}
Too slow for network weights on larger graphs, but feasible for the edge variables on Cora,
where it serves as a reference for the cheaper methods.

\q{Will this scale to big graphs?}
The edge posterior is per edge, so it is linear in the number of edges; the LLM scoring is
the bottleneck, handled by caching or a cheaper scorer.

\q{What if hypotheses fail?}
Each has an informative negative outcome: if H1 fails, one evidence source suffices; if H2
fails, Bayesian GNNs do not help calibration under attack; if H3 fails, embeddings match an
LLM at lower cost.

\q{How does this relate to your research?}
It builds on a research setting I work on (robust message passing on text graphs). The class
project focuses on the Bayesian formulation: edge posteriors, uncertainty and calibration.

%================================================================
\section{Risks and mitigations}
\begin{center}
\small
\begin{tabularx}{\linewidth}{>{\raggedright\arraybackslash}X>{\raggedright\arraybackslash}X}
\toprule
Risk & Mitigation \\
\midrule
LLM scoring too slow or costly & small open model, cache scores, SBERT on large graphs \\
EM loop unstable & damp updates, or compute $d_{ij}$ from a fixed text encoder \\
Mixture likelihood fits poorly & posterior predictive check on $d_{ij}$; try a different family \\
Attack code is complex & start with random and BM25-only injection, add PGD later \\
Results differ by dataset & that is H3's point; report per dataset \\
\bottomrule
\end{tabularx}
\end{center}

%================================================================
\section{Delivery tips and final checklist}
\begin{itemize}
  \item Open by stating the question, not your name and course.
  \item On slide~4, physically point at the three cards in order, then at the bottom strip.
  \item Read each hypothesis's \emph{title} only; let the audience read the details.
  \item If asked something you don't know: ``Good question, I haven't tested that yet; my
        plan would be \ldots''
  \item End on the ask, and pause for questions.
\end{itemize}
\begin{itemize}
  \item[$\square$] Rehearse twice with a timer; aim for 2:45.
  \item[$\square$] Explain the slide~4 equation and the worked example without notes.
  \item[$\square$] Confirm dataset counts and that using the research setting is fine with
        your advisor.
  \item[$\square$] Write the HuskyCT submission (data plus three hypotheses) in your own words.
  \item[$\square$] Bring the deck as both .pptx and PDF.
\end{itemize}

\end{document}
